\section{True Cartier denominators on the original normal curve}

\begin{theorem}[The original curve's actual local rings]
\label{thm:fc-curve-normal-stalk-principal}
\leanok
\lean{EnvelopeProperty.CurveDenominatorCartier.stalk_krullDimLE_one, EnvelopeProperty.CurveDenominatorCartier.stalk_isPrincipalIdealRing}
Let $X$ be the original integral locally Noetherian scheme of topological
dimension at most one, and assume its original local rings are integrally
closed. Every original local ring is a principal ideal ring. The assertion
includes the original generic field and needs no nonfield premise.
\end{theorem}
\begin{proof}
The original topological dimension bounds the coheight of every original
point. Native local-ring dimension equals that coheight, so every original
local ring has Krull dimension at most one. Any nonzero prime in that
local domain is consequently maximal, hence is the actual unique maximal
ideal. The genuine Noetherian local domain equivalences, together with
actual integral closedness, imply that the original local ring is a
principal ideal ring, including the field case.
\end{proof}

\begin{theorem}[Every original rational denominator ideal is Cartier]
\label{thm:fc-curve-denominator-cartier}
\leanok
\lean{EnvelopeProperty.CurveDenominatorCartier.isInvertible_of_nonzero_stalks, EnvelopeProperty.CurveDenominatorCartier.denominator_isInvertible}
\uses{thm:fc-curve-normal-stalk-principal, thm:fc-scheme-rational-denominator-stalks}
On the same original curve, every original ideal with nonzero original
stalks is Cartier. In particular the constructed denominator ideal of
EVERY original $f\in K(X)$ is Cartier, without a supplied nonzero-ideal,
principal-stalk, local-generator or Cartier certificate.
\end{theorem}
\begin{proof}
In each genuine original principal local domain, choose the generator of
the actual original stalk ideal. A nonzero stalk ideal has a nonzero
chosen generator, hence a nonzerodivisor. The genuine locally Noetherian
stalkwise invertibility theorem spreads these generators to original
affine neighborhoods and proves Cartier invertibility. Apply this to
the denominator sheaf, whose same original stalks were already proved
nonzero by actual fraction-field surjectivity.
\end{proof}
